\begin{description}
\item[Q1.] \textbf{Empty concrete parameter window.} Poremba's noise-ratio
constraint $\sqrt{8(m{+}1)}\le 1/\alpha \le q/\sqrt{8(m{+}1)}$ is
\emph{empty} at the wave-brief's small parameters ($n{=}8,q{=}257,m{=}128$)
because $q < 8(m{+}1)$. Correctness still holds empirically at
$\sigma{=}3.2$; but what is the smallest asymptotic instance where the
window is non-empty and simulation is still feasible for verifying
IND-CPA-style attacks? A follow-on study should sweep $(n, q, m)$ and
report both the theoretical lower-bound curve and the largest instance
that fits on a workstation.
\item[Q2.] \textbf{Symmetry between Fourier and computational Del.}
Construction~1 uses \emph{Fourier}-basis $\mathsf{Del}$ and
\emph{computational}-basis $\mathsf{Dec}$, whereas the BB84 primitive it
generalizes uses the opposite convention (Hadamard for cert, computational
for decrypt). At $q > 2$ the Fourier basis is $q$-dimensional and not just
$\{|+\rangle,|-\rangle\}$; the deletion certificate $\pi \in \mathbb{Z}_q^{m+1}$
carries $\log_2 q$ bits per coordinate, not $1$. Does the extra bandwidth
change the cheater's optimal-strategy tradeoff (is it still exponential),
or does it open new sub-exponential attacks that a qubit-only simulation
would miss?
\item[Q3.] \textbf{Strong Gaussian-collapsing conjecture (assumption 1).}
Section~5.2 introduces the ``strong Gaussian-collapsing'' property of
Ajtai's hash function as a \emph{new} assumption --- weaker than
sub-exponential LWE, but not known to reduce to it. Our simulation cannot
falsify it, but numerical experiments \emph{could} search for empirical
distinguishing advantages of quantum adversaries against Ajtai-hash
Gaussian collapsing at moderate $(n,q,m)$. What is the smallest instance
where such a search is meaningful (say, adversary advantage $\gtrsim 0.01$
at $10^6$ samples)?
\item[Q4.] \textbf{Which coordinates of the LWE ciphertext carry deletable
information?} In our BB84 pipeline we picked the \emph{first} 17 of $m{+}1=129$
coordinates and used a simple $b_i = \mathbb{1}[c_i \in (q/4, 3q/4)]$
thresholding. The paper's construction encodes the message only in the
last coordinate (the $b\lfloor q/2\rfloor$ shift). What if we pick the
last 17 coordinates instead --- does the deletion cert on that region
enjoy stronger cheater-independence bounds? A sweep across
coordinate-selection strategies would tell us whether ``deletion-carrying''
coordinates are location-dependent or uniform along the ciphertext.
\item[Q5.] \textbf{Extension to Poremba's actual FHE scheme (\S9).} The
paper's headline application is (leveled) FHE with certified deletion,
which lifts the Dual-Regev primitive through a Mahadev-style FHE compiler.
Simulating the FHE layer at small parameters (bootstrapping omitted, a
few levels of add/multiply) would let us empirically check whether the
certified-deletion accept probability degrades multiplicatively with the
number of homomorphic operations, or remains constant. Neither the paper
nor the [BI20] predecessor exhibits an empirical multi-op experiment;
this is a genuine gap that a follow-on replication could close in maybe
1--2 days of engineering effort.
\end{description}
